% GENERATED FILE. Edit canonical lessons, then run npm run generate:books.
% canonical-source-hash: fnv1a64:a887b93c45dda18c
% canonical-lessons: ZH-C536-chengzhen, ZH-C536-yuanmeng, ZH-C536-meihao, ZH-C536-mengjian, ZH-C536-pantou, ZH-R536-review-532-536

\chapter{When Wishes Come True}
\label{ch:zh-when-wishes-come-true}
\hlchaptermodality{536}

\begin{chapteropening}
\textbf{By the end of this chapter:} \emph{I can say a wish comes true, a dream is fulfilled and life gets better, tell a dream, and say there's something to look forward to.}

Complete the last lesson of chapter 536: I can say a wish comes true, a dream is fulfilled and life gets better, tell a dream, and say there's something to look forward to.
\end{chapteropening}

\section[\texorpdfstring{\zh{成真} (chéngzhēn)}{成真 (chéngzhēn)}]{\zh{成真} (chéngzhēn) — to become reality (of a wish)}

\label{lesson:ZH-C536-chengzhen}



\begin{quote}
Before the new one: say the Mandarin for of one's own free will, then the Mandarin for let us hope that.
\end{quote}

\subsection*{You'll want to know: \zh{成真}}
\textbf{\zh{成真}} — \emph{chéngzhēn} — \textquotedblleft{}to become reality (of a wish)\textquotedblright{}.

\begin{quote}
\textbf{\zh{希望我的愿望能成真}} \emph{xīwàng wǒ de yuànwàng néng chéngzhēn} — I hope my wish comes true.
\end{quote}

\begin{cousinweb}
\textbf{\zh{成真}} (\emph{chéngzhēn}) is built from \zh{成} 'to become' and \zh{真} 'true', both inherited from Middle Chinese, \emph{dzyeng} and \emph{tsyin} (Guangyun, in Baxter's transcription), which glosses the second as the true against the false.
\end{cousinweb}

\begin{grammarlens}[title={What you've built}]
One more word for saying what would happen.
\end{grammarlens}

\subsection*{Your turn}
\begin{itemize}
\raggedright
  \item \emph{Say it:} \emph{chéngzhēn}
  \item \emph{Say it:} \emph{chéngzhēn}, once more
  \item \emph{Say it:} \emph{xīwàng wǒ de yuànwàng néng chéngzhēn} — I hope my wish comes true.
  \item \emph{Recall:} say the Mandarin for of one's own free will, then the Mandarin for let us hope that, then say \emph{chéngzhēn} again
\end{itemize}

\begin{tcolorbox}[breakable,colback=teal!4,colframe=teal!35!black,title={Before you move on}]
What does \zh{成真} mean? (\textquotedblleft{}To become reality (of a wish)\textquotedblright{}.) Say it once more.
\end{tcolorbox}
\section[yuánmèng]{yuánmèng — to fulfil a dream}

\label{lesson:ZH-C536-yuanmeng}



\begin{quote}
Before the new one: say the Mandarin for let us hope that, then the Mandarin for to become reality.
\end{quote}

\subsection*{You'll want to know: yuánmèng}
\emph{yuánmèng} — \textquotedblleft{}to fulfil a dream\textquotedblright{}.

\begin{quote}
\emph{rúguǒ míngnián néng qù Zhōngguó wǒ jiù yuánmèng le} — If I can go to China next year, my dream will have come true.
\end{quote}

Its characters come later; for now it is a word for the ear.

\begin{cousinweb}
\emph{yuánmèng} is built from \emph{yuán} 'round, whole' and \emph{mèng} 'a dream', both inherited from Middle Chinese, \emph{hjwen} and \emph{mjuwngH} (Guangyun, in Baxter's transcription), which glosses the second, after an older dictionary, as waking within sleep: a dream made whole.
\end{cousinweb}

\begin{grammarlens}[title={What you've built}]
One more word for saying what would happen.
\end{grammarlens}

\subsection*{Your turn}
\begin{itemize}
\raggedright
  \item \emph{Say it:} \emph{yuánmèng}
  \item \emph{Say it:} \emph{yuánmèng}, once more
  \item \emph{Say it:} \emph{rúguǒ míngnián néng qù Zhōngguó wǒ jiù yuánmèng le} — If I can go to China next year, my dream will have come true.
  \item \emph{Recall:} say the Mandarin for let us hope that, then the Mandarin for to become reality, then say \emph{yuánmèng} again
\end{itemize}

\begin{tcolorbox}[breakable,colback=teal!4,colframe=teal!35!black,title={Before you move on}]
What does yuánmèng mean? (\textquotedblleft{}To fulfil a dream\textquotedblright{}.) Say it once more.
\end{tcolorbox}
\section[měihǎo]{měihǎo — fine and happy (of a life or a future)}

\label{lesson:ZH-C536-meihao}



\begin{quote}
Before the new one: say the Mandarin for to become reality, then the Mandarin for to fulfil a dream.
\end{quote}

\subsection*{You'll want to know: měihǎo}
\emph{měihǎo} — \textquotedblleft{}fine and happy (of a life or a future)\textquotedblright{}.

\begin{quote}
\emph{zhǐyào dàjiā nǔlì shēnghuó jiù huì gèng měihǎo} — As long as everyone works hard, life will get better.
\end{quote}

Its characters come later; for now it is a word for the ear.

Two third tones meet in this word: měi is said with a rising tone before hǎo, and the spelling keeps both marks.

\begin{cousinweb}
\emph{měihǎo} is built from \emph{měi} 'beautiful' and \zh{好} 'good', both inherited from Middle Chinese, \emph{mijX} and \emph{xawX} (Guangyun, in Baxter's transcription), which glosses the first 'good-looking' and, after an older dictionary, 'sweet', and the second 'good, beautiful'.
\end{cousinweb}

\begin{grammarlens}[title={What you've built}]
One more word for saying what would happen.
\end{grammarlens}

\subsection*{Your turn}
\begin{itemize}
\raggedright
  \item \emph{Say it:} \emph{měihǎo}
  \item \emph{Say it:} \emph{měihǎo}, once more
  \item \emph{Say it:} \emph{zhǐyào dàjiā nǔlì shēnghuó jiù huì gèng měihǎo} — As long as everyone works hard, life will get better.
  \item \emph{Recall:} say the Mandarin for to become reality, then the Mandarin for to fulfil a dream, then say \emph{měihǎo} again
\end{itemize}

\begin{tcolorbox}[breakable,colback=teal!4,colframe=teal!35!black,title={Before you move on}]
What does měihǎo mean? (\textquotedblleft{}Fine and happy (of a life or a future)\textquotedblright{}.) Say it once more.
\end{tcolorbox}
\section[mèngjiàn]{mèngjiàn — to dream of; to see in a dream}

\label{lesson:ZH-C536-mengjian}



\begin{quote}
Before the new one: say the Mandarin for to fulfil a dream, then the Mandarin for fine and happy.
\end{quote}

\subsection*{You'll want to know: mèngjiàn}
\emph{mèngjiàn} — \textquotedblleft{}to dream of; to see in a dream\textquotedblright{}.

\begin{quote}
\emph{wǒ zuótiān wǎnshang mèngjiàn wǒ zhòngjiǎng le} — Last night I dreamt I had won a prize.
\end{quote}

Its characters come later; for now it is a word for the ear.

\begin{cousinweb}
\emph{mèngjiàn} is built from \emph{mèng} 'a dream' and \zh{见} 'to see', both inherited from Middle Chinese, \emph{mjuwngH} and \emph{kenH} (Guangyun, in Baxter's transcription), which glosses the second 'to look'.
\end{cousinweb}

\begin{grammarlens}[title={What you've built}]
One more word for saying what would happen.
\end{grammarlens}

\subsection*{Your turn}
\begin{itemize}
\raggedright
  \item \emph{Say it:} \emph{mèngjiàn}
  \item \emph{Say it:} \emph{mèngjiàn}, once more
  \item \emph{Say it:} \emph{wǒ zuótiān wǎnshang mèngjiàn wǒ zhòngjiǎng le} — Last night I dreamt I had won a prize.
  \item \emph{Recall:} say the Mandarin for to fulfil a dream, then the Mandarin for fine and happy, then say \emph{mèngjiàn} again
\end{itemize}

\begin{tcolorbox}[breakable,colback=teal!4,colframe=teal!35!black,title={Before you move on}]
What does mèngjiàn mean? (\textquotedblleft{}To dream of; to see in a dream\textquotedblright{}.) Say it once more.
\end{tcolorbox}
\section[pàntou]{pàntou — something to look forward to}

\label{lesson:ZH-C536-pantou}



\begin{quote}
Before the new one: say the Mandarin for fine and happy, then the Mandarin for to see in a dream.
\end{quote}

\subsection*{You'll want to know: pàntou}
\emph{pàntou} — \textquotedblleft{}something to look forward to\textquotedblright{}.

\begin{quote}
\emph{yǒu le gōngzuò shēnghuó jiù yǒu pàntou le} — With a job, life has something to look forward to.
\end{quote}

Its characters come later; for now it is a word for the ear.

\begin{cousinweb}
\emph{pàntou} is built from \emph{pàn} 'to long for' and \zh{头}, both inherited from Middle Chinese, \emph{pheanH} and \emph{duw} (Guangyun, in Baxter's transcription). The dictionary glosses the first as beautiful eyes; the sense 'to long for' is not given there. The second, 'the head', is said light here, as at the end of other nouns.
\end{cousinweb}

\begin{grammarlens}[title={What you've built}]
One more word for saying what would happen.
\end{grammarlens}

\subsection*{Your turn}
\begin{itemize}
\raggedright
  \item \emph{Say it:} \emph{pàntou}
  \item \emph{Say it:} \emph{pàntou}, once more
  \item \emph{Say it:} \emph{yǒu le gōngzuò shēnghuó jiù yǒu pàntou le} — With a job, life has something to look forward to.
  \item \emph{Recall:} say the Mandarin for fine and happy, then the Mandarin for to see in a dream, then say \emph{pàntou} again
\end{itemize}

\begin{tcolorbox}[breakable,colback=teal!4,colframe=teal!35!black,title={Before you move on}]
What does pàntou mean? (\textquotedblleft{}Something to look forward to\textquotedblright{}.) Say it once more.
\end{tcolorbox}
\section[(dialogue)]{Looking back: from Hoping for More to When Wishes Come True}

\label{lesson:ZH-R536-review-532-536}



\begin{quote}
Say the Mandarin for with good hopes of it and to be dying to.
\end{quote}

\subsection*{Your turn}
\begin{itemize}
\raggedright
  \item \emph{Say it:} with good hopes of it, to be dying to, to expect much of someone, as one would wish it, to count on someone
  \item \emph{Say it:} to yearn for, to hope that, to thirst for, to give up all hope, may your wishes all come true
  \item \emph{Say it:} would rather, as far as one can, supposing that, to give one's word, may all go as you wish
  \item \emph{Say it:} to be willing, a wish one has, to get what one wished for, of one's own free will, let us hope that
  \item \emph{Say it:} to become reality, to fulfil a dream, fine and happy, to see in a dream, something to look forward to
  \item \emph{Read it:} \zh{望} \zh{希} \zh{愿} — each one aloud, then a word that holds it
\end{itemize}

\begin{tcolorbox}[breakable,colback=teal!4,colframe=teal!35!black,title={Before you move on}]
With good hopes of it? (\textbf{\zh{有望}}, \emph{yǒuwàng}.) To be dying to? (\emph{bābude}.) To expect much of someone? (\textbf{\zh{期望}}, \emph{qīwàng}.) As one would wish it? (\textbf{\zh{如意}}, \emph{rúyì}.) To count on someone? (\emph{zhǐwàng}.) To yearn for? (\emph{xiàngwǎng}.) To hope that? (\textbf{\zh{希望}}, \emph{xīwàng}.) To thirst for? (\emph{kěwàng}.) To give up all hope? (\emph{juéwàng}.) May your wishes all come true? (\textbf{\zh{心想事成}}, \emph{xīnxiǎng shìchéng}.) Would rather? (\emph{nìngkě}.) As far as one can? (\emph{jǐnliàng}.) Supposing that? (\textbf{\zh{假如}}, \emph{jiǎrú}.) To give one's word? (\emph{xǔnuò}.) May all go as you wish? (\textbf{\zh{万事如意}}, \emph{wànshì rúyì}.) To be willing? (\textbf{\zh{愿意}}, \emph{yuànyì}.) A wish one has? (\textbf{\zh{愿望}}, \emph{yuànwàng}.) To get what one wished for? (\textbf{\zh{如愿}}, \emph{rúyuàn}.) Of one's own free will? (\textbf{\zh{自愿}}, \emph{zìyuàn}.) Let us hope that? (\textbf{\zh{但愿}}, \emph{dànyuàn}.) To become reality? (\textbf{\zh{成真}}, \emph{chéngzhēn}.) To fulfil a dream? (\emph{yuánmèng}.) Fine and happy? (\emph{měihǎo}.) To see in a dream? (\emph{mèngjiàn}.) Something to look forward to? (\emph{pàntou}.)
\end{tcolorbox}
